\section{Incertidumbre y predicción selectiva}
\label{sec:marco_incertidumbre}

Decidir cuándo un modelo no debe emitir una predicción requiere responder cuatro preguntas: por qué un modelo
puede equivocarse con la misma seguridad con que acierta, cómo se formaliza la decisión de no responder, con qué
señal se estima la confianza en una predicción y qué puede hacer un sistema, además de responder o no hacerlo,
cuando esa confianza no alcanza.

\subsection{Incertidumbre en la clasificación diagnóstica}
\label{sec:marco_incert_fuentes}

La incertidumbre de una predicción tiene dos fuentes, que se distinguen según si reunir más datos la reduce o
no. La \textbf{incertidumbre aleatoria} proviene de los datos mismos. Una imagen desenfocada, o una lesión cuyo
aspecto dermatoscópico es compatible tanto con un nevo como con un melanoma, deja la respuesta indeterminada
aunque el modelo esté bien entrenado, y más ejemplos de entrenamiento no van a resolver esa incertidumbre. La \textbf{incertidumbre
epistémica} proviene, en cambio, del modelo. Una clase con pocos ejemplos de entrenamiento, como las
minoritarias de un conjunto desbalanceado (Sección~\ref{sec:marco_desbalance}), queda peor representada en lo
que el modelo aprende, y esa incertidumbre disminuye a medida que se agregan datos~\cite{kendall2017uncertainties}.
El cambio de distribución de la Sección~\ref{sec:marco_regularizacion} también produce incertidumbre epistémica,
porque la imagen proviene de un equipo, un centro o una población que el entrenamiento no cubrió.

Ninguna de las dos fuentes queda a la vista en la salida de un clasificador. La función softmax de la
Ecuación~\ref{eq:marco_softmax} reparte siempre entre las clases conocidas una probabilidad que suma uno, de modo
que no existe una salida que exprese «este caso no se parece a nada de lo visto». Una imagen de un tipo de lesión
que el modelo nunca vio recibe igualmente una clase, y la probabilidad de esa clase puede estar cerca de
uno~\cite{hendrycks2017baseline}.

En el caso concreto de dermatología, además, los errores no tienen el mismo costo. No detectar una lesión
maligna tiene consecuencias más graves que derivar una benigna a un especialista, de modo que, ante un caso
dudoso, derivarlo puede ser preferible a arriesgar un diagnóstico. Esa opción existe porque el sistema no toma la
decisión final, sino que asiste a la persona que la toma. En un experimento con dermatólogos y médicos de
atención primaria, el apoyo de un sistema de aprendizaje profundo mejoró la exactitud diagnóstica de los dos
grupos~\cite{groh2024learningaided}. Un sistema que reconoce su límite y deriva el caso puede ser, por lo tanto,
más útil que uno que siempre responde, y medir cuánto se gana al dejar casos sin responder requiere formalizar
esa decisión, que se conoce como \textbf{abstención} (\emph{abstention}).

\subsection{Predicción selectiva y abstención}
\label{sec:marco_pred_selectiva}

Un clasificador con opción de rechazo, o clasificador selectivo, se describe con un par $(f, g)$. La función $f$
asigna una categoría a cada imagen $x$, y la función $g$ decide si la predicción se acepta ($g(x) = 1$) o se
rechaza y el modelo no responde ($g(x) = 0$). La \textbf{cobertura} es la fracción de imágenes que el modelo acepta, y el
\textbf{riesgo selectivo} es la tasa de error medida solo sobre las aceptadas,
\begin{equation}
    \mathrm{cobertura}(g) = \mathbb{E}\bigl[g(x)\bigr], \qquad
    \mathrm{riesgo}(f, g) = \frac{\mathbb{E}\bigl[\mathbf{1}[f(x) \neq y]\; g(x)\bigr]}{\mathbb{E}\bigl[g(x)\bigr]},
    \label{eq:marco_riesgo_selectivo}
\end{equation}
donde $y$ es la categoría real. Sobre $n$ imágenes, la cobertura es la proporción de imágenes con
$g(x_i) = 1$ y el riesgo, la proporción de errores entre ellas~\cite{elyaniv2010foundations,geifman2017selective}.

Una forma habitual de construir $g$ es partir de una señal de confianza $\kappa(x)$, un número que crece con la
seguridad del modelo en su predicción, y fijar un umbral $\tau$,
\begin{equation}
    g(x) = \mathbf{1}\bigl[\kappa(x) \ge \tau\bigr].
    \label{eq:marco_umbral_confianza}
\end{equation}
Subir el umbral reduce la cobertura, porque se aceptan menos imágenes, y reduce también el riesgo selectivo si
la señal es informativa, es decir, si las imágenes que quedan por debajo del umbral son las que el modelo tiende
a clasificar mal. El umbral fija así un compromiso entre cuántos casos se responden y cuántos errores se cometen
entre ellos, y cuánto baja el riesgo por cada caso que se deja de responder depende de la señal elegida.

\subsection{Señales de confianza y calibración}
\label{sec:marco_senales_confianza}

Como la utilidad del umbral depende de la señal, conviene revisar las alternativas para construirla: la
probabilidad de un solo modelo, el desacuerdo entre varios modelos y el voto de los casos más parecidos. La señal
más simple es la máxima probabilidad de la softmax, $\kappa(x) = \max_c p_c$, que sirve de línea base
para detectar predicciones erróneas~\cite{hendrycks2017baseline}. Para que un umbral sobre ella tenga una lectura
probabilística, el modelo tiene que estar \textbf{calibrado}, es decir, sus probabilidades tienen que coincidir
con sus frecuencias de acierto. Por ejemplo, si un modelo calibrado predice melanoma con probabilidad 0.9 en cien
imágenes, alrededor de noventa de ellas son melanomas. En la notación de \textcite{guo2017calibration}, con
$\hat{Y}$ la clase predicha y $\hat{P}$ su probabilidad, la condición es $P(\hat{Y} = Y \mid \hat{P} = p) = p$.
Con un modelo calibrado, un umbral de 0.9 acepta predicciones con al menos 0.9 de probabilidad de ser
correctas; sin calibración, el umbral no tiene esa lectura. \textcite{guo2017calibration} observaron que las
redes neuronales modernas tienden a ser sobreconfiadas, es decir, a asignar a sus predicciones probabilidades
mayores que su frecuencia de acierto, y que las más profundas y anchas suelen estar peor calibradas que las
pequeñas aunque clasifiquen mejor. Las métricas de calibración y las formas de corregirla se presentan en la
Sección~\ref{sec:marco_eval_calib}.

Calibrar la softmax, sin embargo, no hace visible la incertidumbre epistémica. La probabilidad de un solo
modelo, aun calibrada, refleja lo que ese modelo aprendió, y esa incertidumbre corresponde justamente a lo que no
aprendió. Una forma de estimarla es mediante la comparación de modelos. Es esperable
que dos redes entrenadas por separado coincidan donde los datos alcanzaron para fijar la respuesta y difieran
donde no alcanzaron, de modo que su desacuerdo aproxima la incertidumbre epistémica~\cite{kendall2017uncertainties}.
Un \textbf{ensamble} aprovecha este principio. Promedia las probabilidades de varios modelos entrenados por
separado, una combinación que se llama \textbf{voto suave}, y suele ser más exacto y estar mejor calibrado que
cada uno de sus integrantes~\cite{lakshminarayanan2017simple}.

\phantomsection\label{sec:marco_concordancia}%
El desacuerdo puede medirse también entre dos predictores de naturaleza distinta, como un clasificador
supervisado y el \gls{knn} de la Sección~\ref{sec:marco_vlm_clip}. Sus sesgos difieren, porque el primero ajusta
fronteras de decisión entre las clases y el segundo vota con los casos más parecidos a la imagen. Sobre una
misma imagen, sus predicciones pueden coincidir o no: por ejemplo, el clasificador puede predecir un nevo
mientras la mayoría de los vecinos recuperados son melanomas. La \textbf{concordancia} entre sus predicciones
$\hat{y}_1(x)$ y $\hat{y}_2(x)$ es
\begin{equation}
    a(x) = \mathbf{1}\bigl[\hat{y}_1(x) = \hat{y}_2(x)\bigr],
    \label{eq:marco_concordancia}
\end{equation}
y sirve como señal solo si los casos en que los predictores coinciden se equivocan menos que aquellos en que
no, $P(\text{error} \mid a = 1) < P(\text{error} \mid a = 0)$, condición que se verifica con datos. Como usa
las decisiones y no las probabilidades, no depende de que los modelos estén calibrados, pero no detecta los
casos en que los dos predictores cometen el mismo error.

El \textbf{consenso de los vecinos} mide, por su parte, qué tan unánime es el voto de la
Ecuación~\ref{eq:marco_knn}. Es la proporción de vecinos que pertenecen a la clase más votada,
\begin{equation}
    \kappa(x) = \max_{c}\ \hat{p}_{\text{kNN}}(c \mid x),
    \label{eq:marco_consenso}
\end{equation}
que vale 1 cuando los $k$ vecinos pertenecen a la misma clase y disminuye a medida que se reparten entre clases
distintas. Su valor
depende de la representación con que se miden las distancias y de los casos que contiene la base: si la base no
tiene imágenes parecidas a la consultada, los $k$ vecinos se recuperan igual, y su voto no informa sobre ella.

Cualquiera de estas señales se juzga por lo mismo, que es si separa las predicciones correctas de las erróneas.
Una métrica de calibración no mide esa separación, porque la calibración no agrega información que la señal no
contenga. Un modelo que asigna 0.7 a todas sus predicciones y acierta el 70\,\% de ellas está perfectamente
calibrado, pero su confianza no distingue los casos fáciles de los difíciles, y un umbral sobre ella acepta todas
las imágenes o ninguna. Lo que importa de una señal es cuánto menor es la tasa de error entre las predicciones
que acepta que entre las que rechaza. Con una señal que cumple esa condición, además, la decisión puede dejar de
ser binaria.

\subsection{Tres salidas: afirmar, ofrecer diferenciales o derivar}
\label{sec:marco_tres_salidas}

En la práctica clínica, derivar el caso a un especialista no es la única alternativa a afirmar un diagnóstico: un clínico que duda entre
dos o tres diagnósticos los enuncia como \textbf{diagnóstico diferencial}. Un sistema puede tener entonces tres
salidas, que son afirmar un diagnóstico, ofrecer diferenciales, es decir, un conjunto de candidatos, o derivar
el caso a un especialista. Con dos umbrales $\tau_1 > \tau_2$ sobre la señal, la \textbf{regla de decisión} que elige entre
ellas es
\begin{equation}
    \pi(x) =
    \begin{cases}
        \text{afirmar} & \text{si } \kappa(x) \ge \tau_1,\\
        \text{diferencial} & \text{si } \tau_2 \le \kappa(x) < \tau_1,\\
        \text{derivar} & \text{si } \kappa(x) < \tau_2.
    \end{cases}
    \label{eq:marco_politica}
\end{equation}
Si la señal es discreta, como la concordancia o el consenso de unos pocos vecinos, cada salida se asigna a uno o
más de sus valores posibles, y no hay umbrales que ajustar.

La salida diferencial es una predicción por conjuntos, $S(x) \subseteq \{1, \dots, C\}$. El conjunto puede
formarse con las clases más probables de un modelo o con las predicciones de dos predictores que no coinciden, como el par formado por el nevo y el melanoma del ejemplo
anterior. Cuando $|S(x)| = 1$, la salida coincide con afirmar.

Cada salida tiene un costo distinto. Afirmar entrega la información más concreta, pero si la predicción es
errónea, el error llega al lector con la misma forma que un acierto y no puede detectarlo. Ofrecer diferenciales
informa menos, porque deja al lector la elección entre los candidatos, y a cambio es más probable que el
conjunto contenga la clase correcta. Derivar no produce diagnóstico y traslada el caso a una consulta. Cuanto
más costoso es un error frente a una derivación, más alto conviene fijar $\tau_1$, y los umbrales se eligen sobre
el conjunto de validación y no sobre el de test (Sección~\ref{sec:marco_regularizacion}).

Cuando el sistema entrega un reporte clínico, la salida que elige la regla llega al lector a través de un
texto, y un texto puede deshacerla: afirmar con certeza un diagnóstico que la regla dejó abierto entre dos
candidatos, o describir hallazgos que contradicen el diagnóstico que recibió. Por eso, generar el reporte
plantea problemas propios, distintos de los de clasificar.
